\section{Introduction}
\label{sec:introduction}

Green industrial policy is increasingly place based. Governments can support firm-level decarbonizing investment, and policy packages can also include shared inputs---clean-power access, utility networks, ports, industrial land, grid capacity, and related infrastructure---that affect local production conditions. Japan's GX Strategic Area framework illustrates this policy architecture by combining place-based industrial-cluster and infrastructure measures with a separate firm-facing investment-support channel \citep{METIGX2026,METIGXSubsidy2026}. Subnational decarbonization objectives can also be policy relevant; the OECD documents a substantial subnational role in Japan's green transition \citep{OECDJapan2025}. These choices operate in markets where firms compete across jurisdictional boundaries. A policy adopted in one region can therefore change not only emissions and local investment, but also the product-market position of firms located elsewhere.

This raises a cross-instrument public-finance question. Suppose a rival jurisdiction increases firm-facing green support. Should a local government match that support, retreat, or respond through a different instrument? The distinction matters because a rival subsidy may either crowd out local infrastructure provision or induce a jurisdiction to reinforce shared productive capacity. We show that the sign of the infrastructure response can change with product-market substitutability when firms choose both conventional cost-reducing investment and green investment. Throughout, ``switching'' refers to this reversal in the direction of the cross-instrument response, not to discrete adoption of one instrument instead of another.

The contribution is deliberately narrower than the ingredients of the model. Multi-instrument jurisdictional competition and cross-instrument reactions are already central to the fiscal-competition literature; \citet{HauptmeierMittermaierRincke2012} study business taxes and productive public inputs, while \citet{BenassyQuereGobalrajaTrannoy2007} analyze tax and public-input competition more generally. \citet{MoritaOkoshi2025} show that fiscal competition and infrastructure can interact with Cournot competition and threshold effects. Their central comparison concerns how fiscal competition changes infrastructure investment in an FDI setting; our object instead is the sign of the infrastructure best-response derivative to a rival firm-facing subsidy with fixed-location incumbents and endogenous conventional and green investment. Simultaneous conventional and green R\&D is also known \citep{StrandholmEspinolaMunoz2025}. We therefore do not claim novelty for policy multiplicity, infrastructure competition, dual investment, or threshold effects as such. The surviving theorem-level distinction is that, over a certified nonempty global-SPNE region, product substitutability uniquely reverses the sign of the same local infrastructure response to a rival green-investment subsidy; at the canonical remaining primitives, a matched benchmark without conventional investment lacks that reversal.

We study two jurisdictions, each hosting one fixed-location firm. Local governments first choose a firm-specific green-investment subsidy and productive green public infrastructure. Firms then choose conventional cost-reducing investment and green investment, followed by differentiated Cournot competition. Green investment can both reduce emissions and affect productivity, while infrastructure can affect production costs and territorial emissions. The model contains no hard government budget linking the two policies and no firm-level investment constraint linking conventional and green investment.

The main result is a competition-induced sign reversal. On an interior duopoly branch, primitive sufficient conditions reduce the cross-instrument best-response derivative to the sign of a quartic polynomial in $\theta^2$. Bernstein-basis conditions make that polynomial strictly decreasing, yielding a unique threshold $\theta^*$. The global-equilibrium issue is treated separately: with nonnegative policies, investments, and quantities, sufficiently aggressive policies can change market structure. We therefore solve the duopoly, limit-pricing-kink, and monopoly continuations and show that the canonical rational witness supports a true global policy SPNE on a nonempty open neighborhood. At that witness, $\theta^*\simeq0.774$.
Figure~\ref{fig:threshold-response} previews the sign reversal and contrasts it with the matched benchmark without conventional competitive investment.

The economic mechanism is composite. A rival subsidy changes rival green investment and, when $\mu>0$, product-market competitiveness. Quantity changes feed back into both conventional and green investment. The local government's response then reflects consumer-surplus, producer-surplus/business-stealing, territorial-target, and own-policy interaction terms. Product substitutability changes the relative strength of these margins until the switching index crosses zero. In the canonical high-rivalry region the local government reduces its own firm-specific subsidy but raises infrastructure.

A matched benchmark helps isolate the role of conventional investment without turning it into a universal necessity claim. At the canonical environmental and policy primitives, removing conventional competitive investment eliminates the reversal. However, other admissible primitives can generate a no-conventional-investment reversal. The formal conclusion is therefore that conventional investment can create a qualitative best-response change absent in the matched benchmark, not that every reversal requires two investment margins. Likewise, when green investment has no productivity effect, the rival-subsidy product-market transmission vanishes.

The welfare analysis is secondary. Decentralized governments do not internalize the effects of their policies on rival-region consumer surplus, producer surplus, and territorial emissions objectives, so decentralized and coordinated decisions generally differ. In the canonical example coordination raises total model welfare and increases the fiscal weight of infrastructure, but it also raises output and territorial emissions. The result is therefore a second-best regional-welfare comparison, not a theorem that coordination necessarily produces lower emissions.

The institutional mapping remains deliberately limited. The Japanese programs cited above motivate the coexistence of place-based productive infrastructure and firm-facing green investment support, but they do not imply that Japanese local governments literally control both instruments as specified here. The model is a stylized interjurisdictional-policy game rather than a literal mapping of legal authority or financing incidence.

The remainder of the paper proceeds as follows. Section~\ref{sec:model} presents the model and action sets. Section~\ref{sec:equilibrium} solves the interior and boundary continuations. Section~\ref{sec:results} derives the switching threshold and matched benchmark. Section~\ref{sec:welfare} studies coordination. Section~\ref{sec:robustness} considers differentiated Bertrand competition and partial local ownership. Section~\ref{sec:institutional} develops empirical predictions and the institutional mapping. Section~\ref{sec:literature} places the result in the related literature, and the final sections discuss scope and conclude.
